% Standalone source-record documentation; no research attempts.
\documentclass[11pt]{article}
\usepackage[margin=25mm]{geometry}
\usepackage{fontspec,unicode-math}
\setmainfont{Latin Modern Roman}
\setmathfont{Latin Modern Math}
\usepackage{amsmath,amssymb,mathtools,xurl,hyperref}
\hypersetup{colorlinks=true,urlcolor=blue,linkcolor=blue}
\setcounter{secnumdepth}{0}
\setlength{\parindent}{0pt}
\setlength{\parskip}{5pt}
\setlength{\emergencystretch}{3em}
\newcommand{\sref}[2]{\hyperlink{source:#1:#2}{[S#2]}}
\newcommand{\cataloguescope}[1]{\paragraph{Source record.}#1}
\newcommand{\floor}[1]{\left\lfloor#1\right\rfloor}
\begin{document}
\section{Source repair and a sharp quadratic specialization theorem}\label{statement}

\subsection{Definitions and mathematical statement}

\paragraph{Shared notation.}Standard number systems are $\mathbb N,\mathbb Z,\mathbb Q,\mathbb R,\mathbb C$. Specialized notation and index ranges follow the cited source and the definitions below.


Polynomial factorization writes a polynomial as a product of irreducible factors over a specified coefficient field. Absolute factorization allows an algebraic closure. A lattice basis is a finite set of independent integer vectors generating a discrete additive subgroup; basis reduction seeks a geometrically better basis. Approximate factorization instead measures distance between coefficient vectors after a norm is chosen. Structured condition numbers quantify sensitivity under perturbations respecting a matrix's defining structure.
\paragraph{Record scope and qualifications.}
The catalogue title is retained as metadata. Named certificates, contracts, bounds or frameworks in that title are not asserted as results of the original source extract.
The extraction merges most of two PDF pages, including dense and sparse factorization, number-field factors, approximation, and unnamed challenge problems. Broken degree inequalities and missing norms prevent one exact target; these topics remain separate source questions. The supplied catalogue marks this record open; its present research status was not independently checked. The extraction is labelled unrecoverable and is retained as a literal source record; no exact conjecture is reconstructed.
The source text below is rendered literally where necessary; damaged formulas and composite extracts are not silently repaired into a new mathematical claim.
The supplied catalogue marks this record as unrecoverable. Its extract and terminology are documented, but a complete original problem has not been reconstructed.
\paragraph{Statement recorded in the supplied source.}
1. How long does the gradual feeding (or some modification) algorithm actu-ally take to solve one approximate linear equation (with the ci's bounded by kn bits)? (Mark van Hoeij's Question) 2. (a) Can we do lattice basis reduction in essentially linear time in the total number of digits of the input? (Dan Bernstein's Question) (b) Is there a softly linear time reduction from linear system solving to basis reduction in lattices? (Joachim von zur Gathen's question)

3.2 Open Question

Can we use Niederreiter's equation to find a faster algorithm to factor dense polynomials in Q[x]? (Shuhong Gao's Question) The main possible advantage of such a technique would be to use linear algebra over Fp instead of lattice basis reduction.

4 Bivariate

4.1 Open Question

For Dense Bivariate polynomials over Fp with total degree, t: Can we improve on the probabilistic complexity ˜O( d3) + Factorization Complexity of Univariate, degree d, polynomial. Lecerf believes ˜O( dω ) is possible. The bottleneck is solving d2 equations in d variables.

4.2 Open Question

Can we find low-degree ( ≤ n) factors of a polynomial with t terms and degree

d (t
 d2, n
 d)

4.3 Open Question

Find an Absolute Factorization (or test for irreducibility) over fields with small (p < d 2) characteristic, in ˜O( d3).

5 Intermezzo

Given f ∈ Q[x] irreducible, and f (α) = 0. Factor f over Q[α]. State of Art: Belabas 26 Numerical

6.1 Open Question

Challenge Problem \#1 in Kaltofen's JSC 2000 paper.

6.2 Open Question

Given f ∈ C[x, y ] find the nearest polynomial that factors into linear factors in

C[x, y ]. Find/Bound the distance to the nearest polynomial that factors.

6.3 Open Question

Given f ∈ Q[x, y, z ], for which values of z ∈ Q does f factor (completely) in

C[x, y ]? What is the degree of a given z?

6.4 Open Question

How to compute/approximate the structured condition number of a Ruppert Matrix. 3 \sref{statement}{2}
\paragraph{Source status.}
The supplied catalogue labels this source record open; that metadata was checked against the supplied snapshot on 2026-10-11. The label is preserved for provenance and is not an independent verification of present open status or mathematical correctness.
\subsection{Short English statement}
This composite records several questions about lattice reduction, polynomial factorization, and numerical sensitivity.
\subsection{Sources}
\begin{itemize}
\item[\hypertarget{source:statement:1}{S1}] Ulam AI and the UnsolvedMath Contributors, UnsolvedMath: A Curated Collection of Open Mathematics Problems, catalogue entry. Catalogue curation and metadata: CC BY 4.0. The source extract is retained; definitions, English summaries and qualifications are editorial additions.
\url{https://huggingface.co/datasets/ulamai/UnsolvedMath}.
\item[\hypertarget{source:statement:2}{S2}] AIM Workshop Problem Lists, source linked by the catalogue.
\url{https://aimath.org/WWN/polyfactor/polyfactor.pdf}.
\end{itemize}
\paragraph{Preparation and review.}
The original source and UnsolvedMath contributors retain attribution for the source extract. Codex (GPT-6; exact serving revision not exposed) assisted with definitions, English summaries, scope notes and formatting. No new human mathematical authorship or independently verified solution is asserted. Catalogue identity checked against the complete supplied registry and current repository. Distinct upstream records, including repeated questions and broad or damaged extracts, are retained for complete coverage of the source CS category. A repeated record is not a new independent conjecture. Complete mathematical-equivalence and current-literature review has not been performed. Preparation and final review follow contribution prompts C and D, with source gaps disclosed. Structural and rendering checks do not establish a new theorem or open status.
\label{end:statement}
\end{document}
